% Part: first-order-logic
% Chapter: tableaux
% Section: proving-things-quant

\documentclass[../../../include/open-logic-section]{subfiles}

\begin{document}

\olfileid{fol}{tab}{prq}

\olsection{পরিমাণসূচক-সহ \usetoken{P}{tableau}}

\begin{ex}
পরিমাণসূচক নিয়ে কাজ করার সময় আইগেনচল-শর্ত যেন লঙ্ঘিত না হয় তা
নিশ্চিত করতে হবে; কখনও এ জন্য কয়েকটি অনুমান কোন ক্রমে করব তা নিয়ে
পরীক্ষা করতে হয়। সাধারণভাবে, আইগেনচল-শর্তাধীন বিধিগুলি আগে প্রয়োগ
করার চেষ্টা করলে সুবিধা হয় (সম্পূর্ণ !!{tableau}-তে এগুলি ওপরের দিকে
থাকবে)।

$\lexists[x][\lnot !A(x)] \lif \lnot \lforall[x][!A(x)]$
!!{sentence}-টির জন্য কীভাবে !!a{tableau} দেব তা দেখি।
প্রচলিতভাবে সংশ্লিষ্ট অনুমিতি লিখে শুরু করি:
\begin{oltableau}
[\sFmla{\False}{\lexists[x][\lnot \formula{A}(x)]\lif \lnot
    \lforall[x][\formula{A}(x)]}, just=\TAss]
\end{oltableau}
!!{main operator} হলো $\lif$, তাই $\TRule{\False}{\lif}$ প্রয়োগ করি:
\begin{oltableau}
[\sFmla{\False}{\lexists[x][\lnot \formula{A}(x)]\lif \lnot
    \lforall[x][\formula{A}(x)]}, just=\TAss, checked
  [\sFmla{\True}{\lexists[x][\lnot \formula{A}(x)]},
    just={\TRule{\False}{\lif}[1]}
    [\sFmla{\False}{\lnot\lforall[x][\formula{A}(x)]},
      just={\TRule{\False}{\lif}[1]}]
  ]
]
\end{oltableau}
এরপর পংক্তি~$2$ সামলাতে হবে। $\TRule{\True}{\lexists}$ ব্যবহার করি।
এতে একটি নতুন !!{constant} চাই; এখনও কোনো !!{constant}s না থাকায়
যেকোনোটি বেছে নিতে পারি—ধরা যাক~$a$।
\begin{oltableau}
[\sFmla{\False}{\lexists[x][\lnot \formula{A}(x)]\lif \lnot
    \lforall[x][\formula{A}(x)]}, just=\TAss, checked
  [\sFmla{\True}{\lexists[x][\lnot \formula{A}(x)]},
    just={\TRule{\False}{\lif}[1]}, checked
    [\sFmla{\False}{\lnot\lforall[x][\formula{A}(x)]},
      just={\TRule{\False}{\lif}[1]}
      [\sFmla{\True}{\lnot\formula{A}(a)}, just={\TRule{\True}{\lexists}[2]}]
    ]
  ]
]
\end{oltableau}
এবার পংক্তি~$3$-এ $\TRule{\False}{\lnot}$ প্রয়োগ করি:
\begin{oltableau}
[\sFmla{\False}{\lexists[x][\lnot \formula{A}(x)]\lif \lnot
    \lforall[x][\formula{A}(x)]}, just=\TAss, checked
  [\sFmla{\True}{\lexists[x][\lnot \formula{A}(x)]},
    just={\TRule{\False}{\lif}[1]}, checked
    [\sFmla{\False}{\lnot\lforall[x][\formula{A}(x)]},
      just={\TRule{\False}{\lif}[1]}, checked
      [\sFmla{\True}{\lnot\formula{A}(a)}, just={\TRule{\True}{\lexists}[2]}
        [\sFmla{\True}{\lforall[x][\formula{A}(x)]},
          just = {\TRule{\False}{\lnot}[3]}
        ]
      ]
    ]
  ]
]
\end{oltableau}
পংক্তি~$4$-এ $\TRule{\True}{\lnot}$ এবং তারপর পংক্তি~$5$-এ
$\TRule{\True}{\lforall}$ প্রয়োগ করে একটি বদ্ধ !!{tableau} পাই।
\begin{oltableau}
[\sFmla{\False}{\lexists[x][\lnot \formula{A}(x)]\lif \lnot
    \lforall[x][\formula{A}(x)]}, just=\TAss, checked
  [\sFmla{\True}{\lexists[x][\lnot \formula{A}(x)]},
    just={\TRule{\False}{\lif}[1]}, checked
    [\sFmla{\False}{\lnot\lforall[x][\formula{A}(x)]},
      just={\TRule{\False}{\lif}[1]}, checked
      [\sFmla{\True}{\lnot\formula{A}(a)}, just={\TRule{\True}{\lexists}[2]}
        [\sFmla{\True}{\lforall[x][\formula{A}(x)]},
          just = {\TRule{\False}{\lnot}[3]}
          [\sFmla{\False}{\formula{A}(a)}, just={\TRule{\True}{\lnot}[4]}
            [\sFmla{\True}{\formula{A}(a)},
              just={\TRule{\True}{\lforall}[5]}, close
            ]
          ]
        ]
      ]
    ]
  ]
]
\end{oltableau}
\end{ex}

\begin{ex}
নিচের সমষ্টিটির জন্য কীভাবে !!a{tableau} দেব তা দেখি:
\[
\sFmla{\False}{\lexists[x][!C(x,b)]},
\sFmla{\True}{\lexists[x][(!A(x) \land !B(x))]},
\sFmla{\True}{\lforall[x][(!B(x) \lif !C(x,b))]}.
\]
প্রচলিতভাবে অনুমিতিগুলি দিয়ে শুরু করি:
\begin{oltableau}
  [\sFmla{\False}{\lexists[x][\formula{C}(x,b)]}, just=\TAss
    [\sFmla{\True}{\lexists[x][(\formula{A}(x) \land \formula{B}(x))]},
      just=\TAss
      [\sFmla{\True}{\lforall[x][(\formula{B}(x) \lif \formula{C}(x,b))]},
        just=\TAss]
    ]
  ]
\end{oltableau}
আইগেনচল-শর্তযুক্ত বিধি সব সময় আগে প্রয়োগ করা উচিত; এখানে সেটি
পংক্তি~$2$-এ $\TRule{\True}{\lexists}$। অনুমিতিগুলিতে
!!{constant}~$b$ আছে বলে অন্য একটি নিতে হবে; আবার~$a$-ই নিই।
\begin{oltableau}
  [\sFmla{\False}{\lexists[x][\formula{C}(x,b)]}, just=\TAss
    [\sFmla{\True}{\lexists[x][(\formula{A}(x) \land \formula{B}(x))]},
      just=\TAss, checked
      [\sFmla{\True}{\lforall[x][(\formula{B}(x) \lif \formula{C}(x,b))]},
        just=\TAss
        [\sFmla{\True}{\formula{A}(a) \land \formula{B}(a)},
          just={\TRule{\True}{\lexists}[2]}]
      ]
    ]
  ]
\end{oltableau}
এখন পংক্তি~$1$-এ $\TRule{\False}{\lexists}$ অথবা পংক্তি~$3$-এ
$\TRule{\True}{\lforall}$ প্রয়োগ করলে $x$-এর জায়গায় কোন পদ~$t$
প্রতিস্থাপন করব তা ঠিক করতে হয়। এই বিধিগুলিতে আইগেনচল-শর্ত নেই বলে
পছন্দমতো যেকোনো বদ্ধ পদ নেওয়া যায়। কোনো কোনো ক্ষেত্রে ভিন্ন
$t$ দিয়ে বিধিটি কয়েকবারও প্রয়োগ করতে হতে পারে। সাধারণ কৌশল
হিসেবে !!{tableau}-তে আগে থেকেই থাকা কোনো পদ—এখানে $a$ ও~$b$—নিলে
লাভ হয়; এখানে অনুমান করা যায়, $a$ নিলে শাখা বদ্ধ হওয়ার সম্ভাবনা
বেশি।
% BN-SRC-047 TARGET-CORRECTION-BEGIN
\par\smallskip\noindent\textnormal{\small\textbf{উৎস-স্পষ্টীকরণ (BN-SRC-047):} হিমায়িত ইংরেজি গদ্যে আইগেনচল-শর্তহীন দুই বিধিতে “যেকোনো পদ” নেওয়া যায় বলা হয়েছে। পরিমাণসূচক-বিধির ঘোষিত শর্ত অনুযায়ী পদটি বদ্ধ হতে হবে; বাংলা গদ্যে সেই পূর্বশর্তটি স্পষ্ট রাখা হয়েছে।} % BN-SRC-047 TARGET-CORRECTION-END
\begin{oltableau}
  [\sFmla{\False}{\lexists[x][\formula{C}(x,b)]}, just=\TAss
    [\sFmla{\True}{\lexists[x][(\formula{A}(x) \land \formula{B}(x))]},
      just=\TAss, checked
      [\sFmla{\True}{\lforall[x][(\formula{B}(x) \lif \formula{C}(x,b))]}, just=\TAss
        [\sFmla{\True}{\formula{A}(a) \land \formula{B}(a)}, just={\TRule{\True}{\lexists}[2]}
          [\sFmla{\False}{\formula{C}(a,b)}, just={\TRule{\False}{\lexists}[1]}
            [\sFmla{\True}{\formula{B}(a) \lif \formula{C}(a,b)},
              just={\TRule{\True}{\lforall}[3]}
            ]
          ]
        ]
      ]
    ]
  ]
\end{oltableau}
পংক্তি~$1$ ও~$3$-এর !!{signed formula}s-এ টিকচিহ্ন দিই না, কারণ
সেগুলি আবারও ব্যবহার করতে হতে পারে। এবার পংক্তি~$4$-এ
$\TRule{\True}{\land}$ প্রয়োগ করি:
\begin{oltableau}
  [\sFmla{\False}{\lexists[x][\formula{C}(x,b)]}, just=\TAss
    [\sFmla{\True}{\lexists[x][(\formula{A}(x) \land \formula{B}(x))]},
      just=\TAss, checked
      [\sFmla{\True}{\lforall[x][(\formula{B}(x) \lif \formula{C}(x,b))]}, just=\TAss
        [\sFmla{\True}{\formula{A}(a) \land \formula{B}(a)},
          just={\TRule{\True}{\lexists}[2]}, checked
          [\sFmla{\False}{\formula{C}(a,b)}, just={\TRule{\False}{\lexists}[1]}
            [\sFmla{\True}{\formula{B}(a) \lif \formula{C}(a,b)},
              just={\TRule{\True}{\lforall}[3]}
              [\sFmla{\True}{\formula{A}(a)}, just={\TRule{\True}{\land}[4]}
                [\sFmla{\True}{\formula{B}(a)}, just={\TRule{\True}{\land}[4]}
                ]
              ]
            ]
          ]
        ]
      ]
    ]
  ]
\end{oltableau}
এখন পংক্তি~$6$-এ $\TRule{\True}{\lif}$ প্রয়োগ করলে
!!{tableau}-টি বদ্ধ হয়:
\begin{oltableau}
  [\sFmla{\False}{\lexists[x][\formula{C}(x,b)]}, just=\TAss
    [\sFmla{\True}{\lexists[x][(\formula{A}(x) \land \formula{B}(x))]},
      just=\TAss, checked
      [\sFmla{\True}{\lforall[x][(\formula{B}(x) \lif \formula{C}(x,b))]},
        just=\TAss
        [\sFmla{\True}{\formula{A}(a) \land \formula{B}(a)},
          just={\TRule{\True}{\lexists}[2]}, checked
          [\sFmla{\False}{\formula{C}(a,b)},
            just={\TRule{\False}{\lexists}[1]}
            [\sFmla{\True}{\formula{B}(a) \lif \formula{C}(a,b)},
              just={\TRule{\True}{\lforall}[3]}, checked
              [\sFmla{\True}{\formula{A}(a)},
                just={\TRule{\True}{\land}[4]}
                [\sFmla{\True}{\formula{B}(a)},
                  just={\TRule{\True}{\land}[4]}
                  [\sFmla{\False}{\formula{B}(a)},
                    just={\TRule{\True}{\lif}[6]},close
                  ]
                  [\sFmla{\True}{\formula{C}(a,b)},
                    just={\TRule{\True}{\lif}[6]},close
                  ]
                ]
              ]
            ]
          ]
        ]
      ]
    ]
  ]
\end{oltableau}


\end{ex}

\begin{ex}
নিচের সমষ্টিটির জন্য !!a{tableau} তৈরি করি:
\[
\sFmla{\True}{\lforall[x][!A(x)]}, \sFmla{\True}{\lforall[x][!A(x)]
  \lif \lexists[y][!B(y)]}, \sFmla{\True}{\lnot\lexists[y][!B(y)]}.
\]
প্রচলিতভাবে অনুমিতিগুলি লিখে শুরু করি:
\begin{oltableau}
  [\sFmla{\True}{\lforall[x][\formula{A}(x)]}, just=\TAss
    [\sFmla{\True}{\lforall[x][\formula{A}(x)] \lif
        \lexists[y][\formula{B}(y)]}, just=\TAss
      [\sFmla{\True}{\lnot\lexists[y][\formula{B}(y)]}, just=\TAss
      ]
    ]
  ]
\end{oltableau}
পংক্তি~$3$-এ $\TRule{\True}{\lnot}$ প্রয়োগ করে শুরু করি।
“আইগেনচল-শর্তযুক্ত বিধি সব সময় আগে প্রয়োগ করো”—এই নিয়মের একটি
অনুসিদ্ধান্ত হলো “আইগেনচল-শর্তহীন পরিমাণসূচক-বিধির প্রয়োগ প্রয়োজন
হওয়া পর্যন্ত স্থগিত রাখো।” শাখা-বিভাজক বিধিও স্থগিত রাখবে।
\begin{oltableau}
  [\sFmla{\True}{\lforall[x][\formula{A}(x)]}, just=\TAss
    [\sFmla{\True}{\lforall[x][\formula{A}(x)] \lif
        \lexists[y][\formula{B}(y)]}, just=\TAss
      [\sFmla{\True}{\lnot\lexists[y][\formula{B}(y)]}, just=\TAss, checked
        [\sFmla{\False}{\lexists[y][\formula{B}(y)]}, just={\TRule{\True}{\lnot}[3]}]
      ]
    ]
  ]
\end{oltableau}
নতুন পংক্তি~$4$-এ আইগেনচল-শর্তহীন পরিমাণসূচক-বিধি
$\TRule{\False}{\lexists}$ চাই। তাই সেটি স্থগিত রেখে পংক্তি~$2$-এ
$\TRule{\True}{\lif}$ ব্যবহার করি।
\begin{oltableau}
  [\sFmla{\True}{\lforall[x][\formula{A}(x)]}, just=\TAss
    [\sFmla{\True}{\lforall[x][\formula{A}(x)] \lif
        \lexists[y][\formula{B}(y)]}, just=\TAss, checked
      [\sFmla{\True}{\lnot\lexists[y][\formula{B}(y)]}, just=\TAss, checked
        [\sFmla{\False}{\lexists[y][\formula{B}(y)]},
          just={\TRule{\True}{\lnot}[3]},
          [\sFmla{\False}{\lforall[x][\formula{A}(x)]}, just={\TRule{\True}{\lif}[2]}]
          [\sFmla{\True}{\lexists[y][\formula{B}(y)]}, just={\TRule{\True}{\lif}[2]}]
        ]
      ]
    ]
  ]
\end{oltableau}
নতুন দুটি !!{signed formula}s-এই আইগেনচল-শর্তযুক্ত বিধি চাই, তাই
পরের ধাপে সেগুলিই প্রয়োগ করা উচিত:
\begin{oltableau}
  [\sFmla{\True}{\lforall[x][\formula{A}(x)]}, just=\TAss
    [\sFmla{\True}{\lforall[x][\formula{A}(x)] \lif
        \lexists[y][\formula{B}(y)]}, just=\TAss, checked
      [\sFmla{\True}{\lnot\lexists[y][\formula{B}(y)]}, just=\TAss, checked
        [\sFmla{\False}{\lexists[y][\formula{B}(y)]},
          just={\TRule{\True}{\lnot}[3]}
          [\sFmla{\False}{\lforall[x][\formula{A}(x)]}, just={\TRule{\True}{\lif}[2]},checked
            [\sFmla{\False}{\formula{A}(b)}, just={\TRule{\False}{\lforall}[5]}]
          ]
          [\sFmla{\True}{\lexists[y][\formula{B}(y)]}, just={\TRule{\True}{\lif}[2]},checked
            [\sFmla{\True}{\formula{B}(c)}, just={\TRule{\True}{\lexists}[5]}]
          ]
        ]
      ]
    ]
  ]
\end{oltableau}
শাখাগুলি বদ্ধ করতে পংক্তি~$1$ ও~$4$-এর !!{signed formula}s ব্যবহার
করতে হবে। সংশ্লিষ্ট বিধি (\TRule{\True}{\lforall} ও
\TRule{\False}{\lexists})-তে আইগেনচল-শর্ত নেই, তাই উপযুক্ত যেকোনো
বদ্ধ পদ নেওয়া যায়। এখানে যথাক্রমে $b$ ও~$c$ নিতে হবে।
% BN-SRC-048 TARGET-CORRECTION-BEGIN
\par\smallskip\noindent\textnormal{\small\textbf{উৎস-সংশোধন (BN-SRC-048):} হিমায়িত ইংরেজি গদ্যে শেষ ট্যাবলোর জন্য পুনর্ব্যবহারযোগ্য চিহ্নিত সূত্রগুলিকে পংক্তি $1$ ও~$3$ বলা হয়েছে; কিন্তু $\sFmla{\False}{\lexists[y][!B(y)]}$ হলো পংক্তি~$4$, এবং নিচের বৃক্ষও \TRule{\False}{\lexists} বিধির সমর্থন হিসেবে~$[4]$ দেয়। বাংলা গদ্যে পংক্তি~$4$ পুনরুদ্ধার করা হয়েছে।} % BN-SRC-048 TARGET-CORRECTION-END
\begin{oltableau}
  [\sFmla{\True}{\lforall[x][\formula{A}(x)]}, just=\TAss
    [\sFmla{\True}{\lforall[x][\formula{A}(x)] \lif
        \lexists[y][\formula{B}(y)]}, just=\TAss, checked
      [\sFmla{\True}{\lnot\lexists[y][\formula{B}(y)]}, just=\TAss, checked
        [\sFmla{\False}{\lexists[y][\formula{B}(y)]},
          just={\TRule{\True}{\lnot}[3]}
          [\sFmla{\False}{\lforall[x][\formula{A}(x)]},
            just={\TRule{\True}{\lif}[2]},checked
            [\sFmla{\False}{\formula{A}(b)},
              just={\TRule{\False}{\lforall}[5]}
              [\sFmla{\True}{\formula{A}(b)},
                just={\TRule{\True}{\lforall}[1]},close
              ]
            ]
          ]
          [\sFmla{\True}{\lexists[y][\formula{B}(y)]},
            just={\TRule{\True}{\lif}[2]},checked
            [\sFmla{\True}{\formula{B}(c)},
              just={\TRule{\True}{\lexists}[5]}
              [\sFmla{\False}{\formula{B}(c)},
                just={\TRule{\False}{\lexists}[4]},close
              ]
            ]
          ]
        ]
      ]
    ]
  ]
\end{oltableau}
\end{ex}

\begin{prob}
নিচেরগুলির বদ্ধ !!{tableau}s দাও:
\begin{enumerate}
\item $\sFmla{\False}{(\lforall[x][!A(x)] \land \lforall[y][!B(y)])
\lif \lforall[z][(!A(z) \land !B(z))]}$।
\item $\sFmla{\False}{(\lexists[x][!A(x)] \lor \lexists[y][!B(y)])
\lif \lexists[z][(!A(z) \lor !B(z))]}$।
\item $\sFmla{\True}{\lforall[x][(!A(x) \lif !B)]},
\sFmla{\False}{\lexists[y][!A(y)] \lif !B}$।
\item $\sFmla{\True}{\lforall[x][\lnot !A(x)]},
\sFmla{\False}{\lnot\lexists[x][!A(x)]}$।
\item $\sFmla{\False}{\lnot\lexists[x][!A(x)] \lif \lforall[x][\lnot
!A(x)]}$।
\item $\sFmla{\False}{\lnot\lexists[x][\lforall[y][((!A(x,y) \lif
\lnot !A(y,y)) \land (\lnot !A(y,y) \lif !A(x,y)))]]}$।
\end{enumerate}
\end{prob}

\begin{prob}
নিচেরগুলির বদ্ধ !!{tableau}s দাও:
\begin{enumerate}
\item $\sFmla{\False}{\lnot\lforall[x][!A(x)] \lif \lexists[x][\lnot!A(x)]}$।
\item $\sFmla{\True}{(\lforall[x][!A(x)] \lif !B)}, \sFmla{\False}{\lexists[y][(!A(y) \lif !B)]}$।
\item $\sFmla{\False}{\lexists[x][(!A(x) \lif \lforall[y][!A(y)])]}$।
\end{enumerate}
\end{prob}

\end{document}
